\documentclass[spanish,letterpaper,oneside]{article}

\def\documenttitle {Foro de semana 1: Teoria del Estado y Constitucion}
\def\documentsubtitle {Participacion del foro diagnostico de semana 1}
\def\documentsubject {Semana 1 - Teoria del Estado y Constitucion}
\def\documentauthor {Martin Jonathan de la Cruz}
\def\coursename {Teoria del Estado y Constitucion}
\def\coursecode {020}
\def\universityname {Universidad Abierta y a Distancia de Mexico}
\def\universityfaculty {Licenciatura en Derecho}
\def\universitydepartment {Teoria del Estado y Constitucion}
\def\universitydepartmentimage {departamentos/UnADM}
\def\universitydepartmentimagecfg {height=1.57cm}
\def\universitylocation {Roma Norte, Ciudad de Mexico}
\def\authortable {
  \begin{tabular}{ll}
    \textbf{Alumno:} & Martin Jonathan de la Cruz \\
    Matricula: & ES2611202040 \\
    Figura docente: & Ana Lilian Sanchez Chavez \\
    Semestre/Bloque: & 2 / 2 \\
    Estado: & Borrador para revision \\
    \multicolumn{2}{l}{Fecha de realizacion: \today} \\
    \multicolumn{2}{l}{\universitylocation}
  \end{tabular}
}

\input{template}
\usepackage{helvet}
\renewcommand{\familydefault}{\sfdefault}
\usepackage{array}
\usepackage{booktabs}
\usepackage{longtable}
\usepackage{calc}
\usepackage{tikz}
\usetikzlibrary{arrows.meta,positioning,calc}
\setcitestyle{authoryear,open={(},close={)}}
\providecommand{\tightlist}{\setlength{\itemsep}{0pt}\setlength{\parskip}{0pt}}
\providecommand{\pandocbounded}[1]{#1}
\providecommand{\passthrough}[1]{#1}
\setlength{\emergencystretch}{3em}

\usepackage[most]{tcolorbox}
\definecolor{unadmGreenDark}{HTML}{174A3A}
\definecolor{unadmGreen}{HTML}{5F8F3A}
\definecolor{unadmPaper}{HTML}{F6F7F2}
\newtcolorbox{forobox}[1][]{%
  enhanced, breakable,
  colback=unadmPaper, colframe=unadmGreen,
  boxrule=0pt, leftrule=3pt, arc=1.2mm,
  left=4mm, right=4mm, top=3mm, bottom=3mm,
  fonttitle=\bfseries\color{white}, coltitle=white, colbacktitle=unadmGreenDark,
  attach boxed title to top left={xshift=4mm,yshift=-2mm},
  boxed title style={colframe=unadmGreenDark,arc=0.6mm},
  #1
}

\def\coverwatermarkenabled {true}
\def\coverwatermarkimage {img/departamentos/UnADM.pdf}
\def\coverwatermarkopacity {0.16}
\def\coverwatermarkwidth {0.72\paperwidth}
\def\coverwatermarkxshift {0cm}
\def\coverwatermarkyshift {-0.35cm}
\newcommand{\insertcoverwatermark}{%
  \ifthenelse{\equal{\coverwatermarkenabled}{true}}{%
    \AddToShipoutPictureBG*{%
      \begin{tikzpicture}[remember picture,overlay]
        \node[opacity=\coverwatermarkopacity,inner sep=0pt,
          xshift=\coverwatermarkxshift,yshift=\coverwatermarkyshift]
          at (current page.center) {%
            \includegraphics[width=\coverwatermarkwidth]{\coverwatermarkimage}%
          };
      \end{tikzpicture}%
    }%
  }{}%
}

\begin{document}
\insertcoverwatermark
\templatePortrait
\templatePagecfg
\onehalfspacing
\templateIndex
\templateFinalcfg

\section{Identificacion y estado}

Propuesta de participacion basada en el proyecto, elaborada con apoyo de GitHub Copilot. Requiere revision personal; no publicada. No acredita conocimientos previos ni experiencias personales confirmadas. Edad y residencia, cuando se solicitan, permanecen pendientes.

\section{Participacion principal}
\begin{forobox}[title={Participacion principal -- propuesta no publicada}]

\textbf{1. Indique su edad, lugar de residencia actual y su principal
pasatiempo.}

Mi edad es {[}completar edad{]} y actualmente resido en {[}completar
residencia{]}. Entre las actividades que disfruto estan la lectura, la
investigacion y el uso de herramientas digitales para organizar
informacion y desarrollar proyectos academicos.

\textbf{2. ¿Que esperas de la materia?}

Espero comprender como surgio el Estado, como se organiza y que limites
debe respetar frente a las personas. Me interesa relacionar la teoria
con la Constitucion mexicana y analizar por que la existencia de leyes
no siempre garantiza que los derechos se protejan efectivamente. Busco
fortalecer mi capacidad para argumentar y distinguir las funciones de
las instituciones.

\textbf{3. ¿Que es derecho, desde su punto de vista?}

Entiendo el derecho como un conjunto de normas, principios e
instituciones que organiza la convivencia, reconoce derechos y establece
obligaciones. No lo relaciono solamente con sanciones: tambien permite
resolver conflictos y limitar el poder. Su importancia esta en ofrecer
criterios para justificar decisiones y proteger a las personas frente a
actuaciones arbitrarias.

\textbf{4. ¿Enuncia las fuentes formales del derecho?}

Identifico la legislacion, la costumbre y la jurisprudencia. Tambien
considero relevantes los principios generales del derecho y la doctrina,
aunque su funcion y reconocimiento como fuentes dependen del
ordenamiento y de la clasificacion utilizada. Me interesa profundizar en
esas diferencias para no tratar todas las fuentes como si tuvieran la
misma autoridad.

\textbf{5. ¿Que es el Estado de Derecho, desde su punto de vista?}

Es una forma de organizacion en la que tanto las personas como las
autoridades estan sujetas al derecho. Considero que no basta con que
existan leyes: deben protegerse los derechos, limitarse el poder y
existir mecanismos para revisar las decisiones de las autoridades. La
legalidad debe servir para evitar la arbitrariedad, no para darle una
apariencia de justificacion.

\end{forobox}

\section{Respuestas y alcance}

No se agregan respuestas ficticias. La consigna consultada no exige replicas; Antecedentes indica expresamente que no es necesario responder a companeros. Si el docente solicita respuestas adicionales, deben documentarse en cajas independientes tras leer las intervenciones correspondientes.

\end{document}